% GENERATED FILE. Edit canonical lessons, then run npm run generate:books.
% canonical-source-hash: fnv1a64:27f26bf6741bc3fb
% canonical-lessons: ZH-C140-diyi, ZH-C140-yiban, ZH-W140-ban-half, ZH-C140-ci, ZH-C140-zui, ZH-C140-suan, ZH-R140-review-135-140

\chapter{First and Half}
\label{ch:zh-first-and-half}
\hlchaptermodality{140}

\begin{chapteropening}
\textbf{By the end of this chapter:} \emph{I can say first, half, how many times, the most, and work a sum out.}

Complete the last lesson of chapter 140: I can say first, half, how many times, the most, and work a sum out.
\end{chapteropening}

\section[dìyī]{dìyī — first}

\label{lesson:ZH-C140-diyi}



\begin{quote}
Before the new one: say the Mandarin for round, then the Mandarin for pink.
\end{quote}

\subsection*{You'll want to know: dìyī}
\emph{dìyī} — \textquotedblleft{}first\textquotedblright{}.

\begin{quote}
\emph{zhè shì wǒ de dìyī ge shǒujī} — This is my first phone.
\end{quote}

Its characters come later; for now it is a word for the ear.

\zh{一} keeps its first tone here: in a number that gives an order, yī does not change its tone.

\begin{cousinweb}
\emph{dìyī} is built from \emph{dì} 'order, rank' and \zh{一} 'one', both inherited from Middle Chinese, \emph{dejH} and \emph{'jit} (Guangyun, in Baxter's transcription).
\end{cousinweb}

\begin{grammarlens}[title={What you've built}]
One more word for counting.
\end{grammarlens}

\subsection*{Your turn}
\begin{itemize}
\raggedright
  \item \emph{Say it:} \emph{dìyī}
  \item \emph{Say it:} \emph{dìyī}, once more
  \item \emph{Say it:} \emph{zhè shì wǒ de dìyī ge shǒujī} — This is my first phone.
  \item \emph{Recall:} say the Mandarin for round, then the Mandarin for pink, then say \emph{dìyī} again
\end{itemize}

\begin{tcolorbox}[breakable,colback=teal!4,colframe=teal!35!black,title={Before you move on}]
What does dìyī mean? (\textquotedblleft{}First\textquotedblright{}.) Say it once more.
\end{tcolorbox}
\section[\texorpdfstring{\zh{一半} (yíbàn)}{一半 (yíbàn)}]{\zh{一半} (yíbàn) — one half; half (of something)}

\label{lesson:ZH-C140-yiban}



\begin{quote}
Before the new one: say the Mandarin for pink, then the Mandarin for first.
\end{quote}

\subsection*{You'll want to know: \zh{一半}}
\textbf{\zh{一半}} — \emph{yíbàn} — \textquotedblleft{}one half; half (of something)\textquotedblright{}.

\begin{quote}
\textbf{\zh{两千的一半是一千}} \emph{liǎngqiān de yíbàn shì yìqiān} — Half of two thousand is one thousand.
\end{quote}

The new character, in pieces: \zh{半} — two short dots on top, two bars, and a long vertical down through them. The next lesson puts it in your hand.

The tone of \zh{一} changes: yī is said yí before the fourth tone of bàn.

\begin{cousinweb}
\textbf{\zh{一半}} (\emph{yíbàn}) is built from \zh{一} 'one' and \zh{半}, inherited from Middle Chinese \emph{'jit} and \emph{panH} (Guangyun, in Baxter's transcription); the dictionary glosses the second 'a thing divided in the middle', the \emph{bàn} of half past.
\end{cousinweb}

\begin{grammarlens}[title={What you've built}]
One more word for counting.
\end{grammarlens}

\subsection*{Your turn}
\begin{itemize}
\raggedright
  \item \emph{Say it:} \emph{yíbàn}
  \item \emph{Say it:} \emph{yíbàn}, once more
  \item \emph{Say it:} \emph{liǎngqiān de yíbàn shì yìqiān} — Half of two thousand is one thousand.
  \item \emph{Recall:} say the Mandarin for pink, then the Mandarin for first, then say \emph{yíbàn} again
\end{itemize}

\begin{tcolorbox}[breakable,colback=teal!4,colframe=teal!35!black,title={Before you move on}]
What does \zh{一半} mean? (\textquotedblleft{}One half; half (of something)\textquotedblright{}.) Say it once more.
\end{tcolorbox}
\section[\texorpdfstring{\zh{半} (bàn)}{半 (bàn)}]{\zh{半} — two dots, two bars, one vertical}

\label{lesson:ZH-W140-ban-half}



\begin{quote}
The word first, then the character inside it.

\begin{itemize}
\raggedright
  \item \emph{Recall:} say \emph{yíbàn} — \textbf{R1}, one lesson back, before it decays
\end{itemize}

\begin{quote}
\zh{二} \zh{千}
\end{quote}

\emph{Read it:} each one aloud, with nothing said first — \emph{èr}, \emph{qiān}
\end{quote}

\subsection*{Script — two dots, two bars, one vertical}
\hlblockfigure{\detokenize{figures/ZH-W140-ban-half-filmstrip.pdf}}{How \zh{半} is written, stroke by stroke}

\begin{quote}
\zh{半}
\end{quote}

\textbf{One piece, five strokes.}

\textbf{Two short dots} at the top, leaning toward each other: the left one falls to the lower right, the right one to the lower left.

\textbf{Two horizontals}, short over long, the two bars of \zh{二}.

\textbf{A long vertical} from above the shorter bar, straight down through both to the foot.

2 + 2 + 1 = five.

\subsection*{Your turn}
\begin{itemize}
\raggedright
  \item \emph{Write it:} \zh{半}, all five, with the model in front of you
  \item \emph{Cover:} write it again from memory, and count the pieces, not the strokes
\end{itemize}

\begin{tcolorbox}[breakable,colback=teal!4,colframe=teal!35!black,title={Before you move on}]
How many strokes in \zh{半}? (\textbf{Five}.) Which way do the two dots lean? (\textbf{Toward each other}.) What comes last? (\textbf{The long vertical}.)

Source: \href{https://raw.githubusercontent.com/chanind/hanzi-writer-data/68d10a4b21150cae5e1ebbd223eed289cf32d90c/data/\%E5\%8D\%8A.json}{Hanzi Writer Data \zh{半}, pinned stroke paths and medians}.
\end{tcolorbox}
\section[cì]{cì — time (one occurrence)}

\label{lesson:ZH-C140-ci}



\begin{quote}
Before the new one: say the Mandarin for first, then the Mandarin for one half.
\end{quote}

\subsection*{You'll want to know: cì}
\emph{cì} — \textquotedblleft{}time (one occurrence)\textquotedblright{}.

\begin{quote}
\emph{zhè shì wǒ dìyī cì lái Běijīng} — This is my first time in Beijing.
\end{quote}

Its characters come later; for now it is a word for the ear.

\begin{cousinweb}
\emph{cì} is inherited from Middle Chinese \emph{tshijH} (Guangyun, in Baxter's transcription), glossed 'order, sequence' and 'a halt on a journey'.
\end{cousinweb}

\begin{grammarlens}[title={What you've built}]
One more word for counting.
\end{grammarlens}

\subsection*{Your turn}
\begin{itemize}
\raggedright
  \item \emph{Say it:} \emph{cì}
  \item \emph{Say it:} \emph{cì}, once more
  \item \emph{Say it:} \emph{zhè shì wǒ dìyī cì lái Běijīng} — This is my first time in Beijing.
  \item \emph{Recall:} say the Mandarin for first, then the Mandarin for one half, then say \emph{cì} again
\end{itemize}

\begin{tcolorbox}[breakable,colback=teal!4,colframe=teal!35!black,title={Before you move on}]
What does cì mean? (\textquotedblleft{}Time (one occurrence)\textquotedblright{}.) Say it once more.
\end{tcolorbox}
\section[zuì]{zuì — most}

\label{lesson:ZH-C140-zui}



\begin{quote}
Before the new one: say the Mandarin for one half, then the Mandarin for one occurrence.
\end{quote}

\subsection*{You'll want to know: zuì}
\emph{zuì} — \textquotedblleft{}most\textquotedblright{}.

\begin{quote}
\emph{wǒ zuì xǐhuan chūntiān} — Spring is my favourite.
\end{quote}

Its characters come later; for now it is a word for the ear.

\begin{cousinweb}
\emph{zuì} is inherited from Middle Chinese \emph{tswajH} (Guangyun, in Baxter's transcription), glossed 'utmost'.
\end{cousinweb}

\begin{grammarlens}[title={What you've built}]
One more word for counting.
\end{grammarlens}

\subsection*{Your turn}
\begin{itemize}
\raggedright
  \item \emph{Say it:} \emph{zuì}
  \item \emph{Say it:} \emph{zuì}, once more
  \item \emph{Say it:} \emph{wǒ zuì xǐhuan chūntiān} — Spring is my favourite.
  \item \emph{Recall:} say the Mandarin for one half, then the Mandarin for one occurrence, then say \emph{zuì} again
\end{itemize}

\begin{tcolorbox}[breakable,colback=teal!4,colframe=teal!35!black,title={Before you move on}]
What does zuì mean? (\textquotedblleft{}Most\textquotedblright{}.) Say it once more.
\end{tcolorbox}
\section[suàn]{suàn — to calculate; to work out}

\label{lesson:ZH-C140-suan}



\begin{quote}
Before the new one: say the Mandarin for one occurrence, then the Mandarin for most.
\end{quote}

\subsection*{You'll want to know: suàn}
\emph{suàn} — \textquotedblleft{}to calculate; to work out\textquotedblright{}.

\begin{quote}
\emph{wǒ suàn le yígòng shì sānshí kuài} — I've worked it out: thirty yuan in all.
\end{quote}

Its characters come later; for now it is a word for the ear.

\begin{cousinweb}
\emph{suàn} is inherited from Middle Chinese \emph{swanX} (Guangyun, in Baxter's transcription), glossed 'the number of things'. The regular sound change would give a third tone, \emph{suǎn}, not \emph{suàn}; why is not settled.
\end{cousinweb}

\begin{grammarlens}[title={What you've built}]
One more word for counting.
\end{grammarlens}

\subsection*{Your turn}
\begin{itemize}
\raggedright
  \item \emph{Say it:} \emph{suàn}
  \item \emph{Say it:} \emph{suàn}, once more
  \item \emph{Say it:} \emph{wǒ suàn le yígòng shì sānshí kuài} — I've worked it out: thirty yuan in all.
  \item \emph{Recall:} say the Mandarin for one occurrence, then the Mandarin for most, then say \emph{suàn} again
\end{itemize}

\begin{tcolorbox}[breakable,colback=teal!4,colframe=teal!35!black,title={Before you move on}]
What does suàn mean? (\textquotedblleft{}To calculate; to work out\textquotedblright{}.) Say it once more.
\end{tcolorbox}
\section[(dialogue)]{Looking back: from Sport to First and Half}

\label{lesson:ZH-R140-review-135-140}



\begin{quote}
Say the Mandarin for football and basketball.
\end{quote}

\subsection*{Your turn}
\begin{itemize}
\raggedright
  \item \emph{Say it:} football, basketball, to kick, a match, to work out
  \item \emph{Say it:} a sunny day, an overcast day, to be windy, to thunder, pleasantly cool
  \item \emph{Say it:} festival, Chinese New Year, to celebrate the New Year, to have time off, the holidays
  \item \emph{Say it:} each, oneself, other people, every single, the others
  \item \emph{Say it:} green, yellow, grey, round, pink
  \item \emph{Say it:} first, one half, one occurrence, most, to calculate
  \item \emph{Read it:} \zh{节} \zh{每} \zh{半} — each one aloud, then a word that holds it
\end{itemize}

\begin{tcolorbox}[breakable,colback=teal!4,colframe=teal!35!black,title={Before you move on}]
Football? (\emph{zúqiú}.) Basketball? (\emph{lánqiú}.) To kick? (\emph{tī}.) A match? (\emph{bǐsài}.) To work out? (\emph{duànliàn}.) A sunny day? (\emph{qíngtiān}.) An overcast day? (\emph{yīntiān}.) To be windy? (\emph{guāfēng}.) To thunder? (\emph{dǎléi}.) Pleasantly cool? (\emph{liángkuai}.) Festival? (\textbf{\zh{节日}}, \emph{jiérì}.) Chinese New Year? (\textbf{\zh{春节}}, \emph{Chūnjié}.) To celebrate the New Year? (\emph{guònián}.) To have time off? (\emph{fàngjià}.) The holidays? (\emph{jiàqī}.) Each? (\textbf{\zh{每}}, \emph{měi}.) Oneself? (\emph{zìjǐ}.) Other people? (\emph{biérén}.) Every single? (\emph{suǒyǒu}.) The others? (\emph{qítā}.) Green? (\emph{lǜ}.) Yellow? (\emph{huáng}.) Grey? (\emph{huī}.) Round? (\emph{yuán}.) Pink? (\emph{fěnhóng}.) First? (\emph{dìyī}.) One half? (\textbf{\zh{一半}}, \emph{yíbàn}.) One occurrence? (\emph{cì}.) Most? (\emph{zuì}.) To calculate? (\emph{suàn}.)
\end{tcolorbox}
